\section{Three-Tier Validation}

\label{sec:validation}
%% ============================================================================

QuasarSTM validates a speculatively executed transaction in three
tiers, falling through only when the previous tier cannot prove
serializability. The tiers are ordered by cost: Tier~1 touches only
lane clocks, Tier~2 walks the read set against per-key MVCC chains,
Tier~3 invokes a registered semantic reducer.

\subsection{Tier 1: Lane-Clock Fast Validation}

\begin{definition}[Tier-1 validation predicate]
\label{def:tier1}
Let $T$ be a speculatively executed transaction with lane snapshot
$\mathit{LaneSnapshot}(T) = \{(L_i, c_i)\}_{i=1}^{m}$. Tier-1
validation succeeds for $T$ iff
\[
\forall i \in [m]:\quad
\mathit{LaneClock}[L_i].\mathit{version} = c_i.
\]
\end{definition}

\begin{lemma}[Tier-1 sufficiency]
\label{lem:tier1-sufficient}
If $T$ passes Tier-1 validation under
$\preceq_{\mathrm{can}}$, then every read recorded by $T$ observed the
visible version under $\preceq_{\mathrm{can}}$, and every write of $T$
is the next visible version on its key.
\end{lemma}

\begin{proof}[Sketch]
Lane clocks advance only on committed writes to the lane
(\S\ref{sec:lanes}). The snapshot $c_i$ was taken before $T$ began;
the equality at validation implies no committed write to $L_i$ has
intervened. Each storage key $k$ touched by $T$ belongs to exactly
one lane $L(k)$; if $\mathit{LaneClock}[L(k)]$ is unchanged, the
visible version of $k$ under $\preceq_{\mathrm{can}}$ at $T$'s
canonical position is unchanged, so $T$'s observed version is
correct.
\end{proof}

Tier-1 is the design's fast path. On the disjoint-lane workload
(\S\ref{sec:performance} bench A), Tier-1 hit rate exceeds 95\%.

\subsection{Tier 2: Key-Level MVCC Validation}

When a lane clock has advanced, Tier-1 cannot prove the read was
unaffected. Tier-2 walks the read set entry-by-entry against the
canonical MVCC chain.

\begin{definition}[Read record]
\label{def:read-record}
Each read recorded by $T$ is a tuple
$\mathit{StmRead} = (\mathit{tx\_id},\, \mathit{key},\,
\mathit{version\_id},\, \mathit{observed\_write\_ts},\,
\mathit{observed\_lane\_clock})$.
\end{definition}

\begin{definition}[Tier-2 validation predicate]
\label{def:tier2}
Tier-2 validation succeeds for $T$ iff for every read record
$(\cdot, k, v, \cdot, \cdot)$ recorded by $T$, the version $v$ is
the visible version of $k$ under $\preceq_{\mathrm{can}}$ at $T$'s
canonical position (\S\ref{sec:visibility}, Definition~\ref{def:visible}).
\end{definition}

A Tier-2 mismatch sends $T$ either to Tier-3 (if all mismatched writes
are in commutative ops) or directly to the repair queue
(\S\ref{sec:repair}).

\subsection{Tier 3: Semantic Commutativity Validation}

Many writes commute deterministically. Treating them as ordinary
\textsc{Set} operations is correct but conservative: it produces
spurious conflicts whenever two transactions \texttt{Add} to the same
counter. Tier-3 admits these commits semantically.

\begin{definition}[Semantic write op]
\label{def:semantic-op}
A write is tagged with a \emph{semantic op}
$\mathit{StmWriteOp} \in \{\textsc{Set},\, \textsc{Add},\, \textsc{Sub},\,
\textsc{Min},\, \textsc{Max},\, \textsc{Append},\, \textsc{BitOr},\,
\textsc{BalanceDelta},\, \textsc{OrderAppend},\,
\textsc{FeeAccumulate}\}$.
\end{definition}

\begin{definition}[Commutative on lane]
\label{def:commutative}
A pair of semantic ops $(o_1, o_2)$ \emph{commutes on lane $L$}
iff applying $o_1$ then $o_2$ to the lane state yields the same
post-state as $o_2$ then $o_1$, for all valid pre-states.
\end{definition}

The commutativity matrix is fixed by specification:

\begin{table}[ht]
\centering
\begin{tabular}{@{}lc@{}}
\toprule
\textbf{Op pair} & \textbf{Commutes?} \\
\midrule
$\textsc{Append} + \textsc{Append}$  & yes (deterministic by canonical index) \\
$\textsc{Add}    + \textsc{Add}$      & yes (sum reducer) \\
$\textsc{FeeAccumulate}$              & yes (sum reducer) \\
$\textsc{BalanceDelta} + \textsc{BalanceDelta}$ & yes if no overdraft, canonical netting \\
$\textsc{OrderAppend}$                & yes within batch-auction lane \\
$\textsc{Min}/\textsc{Max}/\textsc{BitOr}$ & yes (idempotent reducers) \\
$\textsc{Set} + \text{anything}$      & no \\
\bottomrule
\end{tabular}
\caption{Semantic commutativity matrix. Tier-3 admits a transaction
whose conflicting writes are all in commuting cells.}
\label{tab:commutativity}
\end{table}

\begin{definition}[Tier-3 validation predicate]
\label{def:tier3}
Tier-3 validation succeeds for $T$ iff every Tier-2 mismatch is
attributable to a write whose semantic op commutes (under
Definition~\ref{def:commutative}) with every other transaction's
mismatching write on the same lane, and the lane has a registered
\emph{reducer} that combines the contributions deterministically by
canonical index.
\end{definition}

\paragraph{Reducers.}
A lane reducer is a per-lane function $r_L$ with type
$\mathit{Value} \times \mathit{Op}^* \to \mathit{Value}$ that consumes
the lane's pre-state and the canonical-order sequence of contributing
ops to produce the post-state. Reducers are pure, deterministic, and
gas-metered. They run in the
\texttt{StmSemanticValidate} workgroup (service ID \texttt{0x0432})
ahead of commit.

\paragraph{Why Tier-3 is mandatory for regulated DEX.}
The order-book \texttt{Append}, fee accumulation, balance
delta, and audit append lanes dominate exchange workloads.
Without Tier-3, these become hot-key conflicts that repeatedly
abort. With Tier-3, they collapse to a single deterministic reduction
per lane per round. The performance gap is the difference between a
fast blockchain and a fast regulated DEX.

\subsection{Tier Cascade}

\begin{algorithm}[H]
\caption{Validate transaction $T$}
\label{alg:validate}
\begin{algorithmic}[1]
\State \textbf{if} Tier-1 succeeds for $T$ \textbf{then} \Return
       \textsc{LaneValidated}
\State \textbf{if} Tier-2 succeeds for $T$ \textbf{then} \Return
       \textsc{KeyValidated}
\State \textbf{if} Tier-3 succeeds for $T$ \textbf{then} \Return
       \textsc{SemanticallyValidated}
\State \Return \textsc{RepairPending}
\end{algorithmic}
\end{algorithm}

\begin{remark}[Cost monotonicity]
\label{rem:tier-cost}
Tier-1 inspects $|m|$ lane clocks (typically $m \leq 4$). Tier-2
walks $|R(T)|$ read records (typically tens). Tier-3 dispatches a
reducer on the affected lanes. The cost cascade is exactly inverted
from frequency: most transactions stop at Tier~1, which is the
cheapest. The pathology Block-STM~2.0 suffered---always running the
expensive validator---is eliminated by structure.
\end{remark}


%% ============================================================================
